\documentclass[a4paper,12pt]{article}
\usepackage[T2A]{fontenc}
\usepackage[utf8]{inputenc}
\usepackage[russian]{babel}

\usepackage{amsmath, amssymb, amsthm, mathtools}
\usepackage{helvet}
\renewcommand{\familydefault}{\sfdefault}
\usepackage{icomma}
\usepackage[left=18mm,right=18mm,top=18mm,bottom=20mm]{geometry}
\usepackage{microtype}
\usepackage{tikz}
\usepackage{tkz-euclide}
\usepackage{pgfplots}
\pgfplotsset{compat=1.18}
\usepackage{tabularx, booktabs, longtable}
\usepackage{enumitem}
\usepackage{hyperref}
\usepackage{parskip}
\usepackage{fancyhdr}
\usepackage{setspace}
\usepackage{xcolor}

\definecolor{accent}{HTML}{008080}
\definecolor{darkaccent}{HTML}{004D4D}
\definecolor{textgray}{HTML}{333333}
\definecolor{softgray}{HTML}{F2F6F6}
\definecolor{muted}{HTML}{6B7C7C}

\hypersetup{hidelinks}

\onehalfspacing
\setlength{\parindent}{0pt}
\setlength{\headheight}{14.5pt}

\pagestyle{fancy}
\fancyhf{}
\lhead{\textcolor{muted}{Прикладные задачи}}
\rhead{\textcolor{muted}{ЕГЭ}}
\cfoot{\thepage}

\setlist[itemize]{leftmargin=1.2em,itemsep=0.25em,topsep=0.25em}
\setlist[enumerate]{leftmargin=1.4em,itemsep=0.35em,topsep=0.35em}

\newcommand{\SectionTitle}[1]{%
  \vspace{0.8em}
  {\Large\bfseries\textcolor{darkaccent}{#1}}\par
  \vspace{0.2em}
  {\color{accent}\rule{\linewidth}{0.8pt}}
  \vspace{0.4em}
}

\newcommand{\MiniTitle}[1]{%
  \vspace{0.45em}
  {\bfseries\textcolor{darkaccent}{#1}}\par
}

\newcommand{\Important}[1]{%
  \textcolor{darkaccent}{\bfseries #1}
}

\begin{document}

\begin{titlepage}
\centering
\vspace*{28mm}

{\Large\textcolor{darkaccent}{ЕГЭ по математике}\par}
\vspace{16mm}

{\fontsize{25}{30}\selectfont\bfseries Чек-лист\par}
\vspace{5mm}
{\Large\bfseries Прикладные задачи\par}
\vspace{3mm}
{\large экономика, физика, кинематика и термодинамика\par}

\vspace{22mm}

{\large Лёвин Артём Александрович\par}
\vspace{2mm}
{\large эксклюзивно для Ксении Клыковой\par}

\vspace{18mm}

{\large \today\par}

\vfill

\begin{tikzpicture}[remember picture,overlay]
  \draw[
    accent,
    line width=1.15pt
  ]
  ([xshift=22mm,yshift=28mm]current page.south west)
  .. controls
  ([xshift=62mm,yshift=52mm]current page.south west)
  and
  ([xshift=132mm,yshift=8mm]current page.south west)
  ..
  ([xshift=188mm,yshift=30mm]current page.south west);

  \draw[
    darkaccent,
    line width=0.65pt
  ]
  ([xshift=28mm,yshift=18mm]current page.south west)
  .. controls
  ([xshift=82mm,yshift=34mm]current page.south west)
  and
  ([xshift=128mm,yshift=2mm]current page.south west)
  ..
  ([xshift=178mm,yshift=20mm]current page.south west);
\end{tikzpicture}

\end{titlepage}

\SectionTitle{1. Экономические задачи: выручка и цена}

В задачах на монополию часто требуется найти максимальную или минимальную цену, при которой выполняется условие по выручке.

\MiniTitle{Алгоритм решения}
\begin{enumerate}
\item Выпишите данные: зависимость объёма спроса $Q(P)$ и формулу выручки $R = Q \cdot P$.
\item Составьте \Important{каскад подстановок}: подставьте $Q(P)$ в формулу $R$, чтобы получить уравнение с одной переменной $P$.
\item Слова «не менее» или «не более» в граничных задачах заменяем на знак равенства (решаем уравнение $R = R_{\text{граничное}}$).
\item Найдите корни уравнения и выберите тот, который удовлетворяет условию (например, \textit{максимальную} цену).
\end{enumerate}

\MiniTitle{Пример}
Зависимость объёма спроса: $Q = 100 - 5P$. Выручка $R = Q \cdot P$. Найти максимальную цену $P$, при которой $R \ge 375$ (тыс. руб.).

\textit{Решение:}
\[
R = (100 - 5P)P = 100P - 5P^2.
\]
Решаем уравнение $100P - 5P^2 = 375$:
\[
5P^2 - 100P + 375 = 0 \quad \Big| :5 \implies P^2 - 20P + 75 = 0.
\]
По теореме Виета: $P_1 = 5$, $P_2 = 15$.
Так как нужна \textit{максимальная} цена, в ответ записываем $15$.

\SectionTitle{2. Физические задачи и степени десятки}

В задачах на давление, скорость или другие физические величины часто встречаются десятичные дроби и «тысячные».

\MiniTitle{Правило работы со степенями}
Не работайте с десятичными дробями напрямую. Переводите их в стандартный вид (степени десятки).
\[
0{,}625 \text{ тысячных} = 0{,}625 \cdot 10^{-3} = 625 \cdot 10^{-6}.
\]

\MiniTitle{Извлечение корня из степени}
При извлечении корня степени делятся на показатель корня:
\[
\sqrt{625 \cdot 10^{-4}} = \sqrt{625} \cdot \sqrt{10^{-4}} = 25 \cdot 10^{-2} = 0{,}25.
\]

\Important{Типичная ошибка:} забыть, что при сдвиге запятой вправо на $n$ знаков, показатель степени десятки \textit{уменьшается} на $n$.

\SectionTitle{3. Кинематика: время нахождения выше отметки}

Если тело брошено вертикально (или под углом), его высота описывается квадратичной функцией:
\[
h(t) = h_0 + v_0 t - \frac{g}{2} t^2.
\]

\MiniTitle{Смысл корней и отбор}
Уравнение $h(t) = H$ имеет два корня $t_1$ и $t_2$ ($t_1 < t_2$).
\begin{itemize}
\item $t_1$ — время, когда тело достигает высоты $H$ при подъёме.
\item $t_2$ — время, когда тело достигает высоты $H$ при падении.
\end{itemize}

Чтобы найти время, в течение которого тело находилось \textit{выше} отметки $H$, нужно найти разность:
\[
\Delta t = t_2 - t_1.
\]

\begin{center}
\begin{tikzpicture}[scale=0.7]
  \draw[->, line width=0.8pt] (0,0) -- (0,10) node[left]{$h$};
  \draw[->, line width=0.8pt] (0,0) -- (7,0) node[below]{$t$};
  
  \draw[darkaccent, thick, domain=0:6, smooth, variable=\x] plot ({\x}, {6*\x - \x*\x});
  
  \draw[dashed, textgray] (0,5) -- (6.5,5) node[right, textgray]{$H$};
  
  \filldraw[accent] (1,5) circle (2pt) node[below left]{$t_1$};
  \filldraw[accent] (5,5) circle (2pt) node[below right]{$t_2$};
  
  \draw[accent, thick, <->] (1, 5.6) -- (5, 5.6);
  \node[above, accent, font=\small] at (3, 5.6) {$\Delta t = t_2 - t_1$};
\end{tikzpicture}
\end{center}

\SectionTitle{4. Термодинамика: нагрев и разрушение}

Температура нагревательного элемента часто задаётся формулой:
\[
T(t) = T_0 + bt + at^2.
\]

\MiniTitle{Физический смысл корней}
Если прибор портится при достижении температуры $T_{\text{крит}}$, мы решаем уравнение $T(t) = T_{\text{крит}}$.
Квадратное уравнение даст два корня: $t_1$ и $t_2$.

\Important{Главное правило отбора:} прибор достигает критической температуры в \textit{меньший} момент времени ($t_1$). В этот момент он либо ломается, либо отключается. До второго корня ($t_2$) он просто не доживёт, поэтому $t_2$ — математически верный, но физически бессмысленный корень.

\MiniTitle{Пример}
$T(t) = 1400 + 200t - 10t^2$. Прибор портится при $1760$ К.
\[
1400 + 200t - 10t^2 = 1760 \implies 10t^2 - 200t + 360 = 0 \implies t^2 - 20t + 36 = 0.
\]
Корни: $t_1 = 2$, $t_2 = 18$.
Прибор нагреется до $1760$ К за $2$ секунды и отключится. Ответ: $2$.

\SectionTitle{5. Общий алгоритм и ловушки}

\begin{enumerate}
\item \textbf{Игнорируйте «не более/не менее»} при составлении граничного уравнения (заменяйте на $=$).
\item \textbf{Каскад подстановок} в экономике: всегда стремитесь получить уравнение с одной переменной.
\item \textbf{Степени десятки}: переводите $0{,}xxx$ в $A \cdot 10^{-n}$ перед извлечением корня.
\item \textbf{Физический смысл}: в задачах на время или разрушение всегда выбирайте \textit{первый} момент достижения критического состояния.
\item \textbf{Интервалы}: если спрашивают «сколько времени длился процесс», вычисляйте $\Delta t = t_2 - t_1$.
\end{enumerate}

\newpage

\SectionTitle{Тренировочный блок}

Решите задачи. Ответ запишите в виде числа или десятичной дроби.

\begin{enumerate}[label=\textbf{\arabic*.}]
\item 
Зависимость объёма спроса $Q$ (тыс. руб.) от цены $P$ (тыс. руб.) задана формулой $Q = 90 - 5P$. Выручка $R = Q \cdot P$. Определите максимальный уровень цены $P$, при котором месячная выручка $R$ составит не менее $360$ тыс. руб.

\item 
Температура нагревательного элемента изменяется по закону $T(t) = 1000 + 100t - 5t^2$ (в Кельвинах). Прибор необходимо отключить, если температура превышает $1320$ К. Найдите максимальное время (в секундах) работы прибора до отключения.

\item 
Высота мяча над землёй изменяется по закону $h(t) = 10 + 15t - 5t^2$ (в метрах). Определите время (в секундах), в течение которого мяч находится на высоте не менее $20$ метров.

\item 
Скорость воды в трубе определяется формулой $v = \sqrt{0{,}000036}$. Найдите скорость $v$ в м/с, переведя подкоренное выражение в стандартный вид (степени десятки).

\item 
Зависимость объёма спроса $Q$ от цены $P$ задана формулой $Q = 110 - 5P$. Определите максимальный уровень цены $P$, при котором выручка $R$ составит не менее $480$.

\item 
Температура прибора задаётся законом $T(t) = 1200 + 140t - 5t^2$. Прибор необходимо отключить при достижении $1680$ К. Найдите максимальное время работы прибора.

\item 
Камень брошен вверх. Его высота изменяется по закону $h(t) = 1{,}2 + 17t - 5t^2$. Найдите время, в течение которого камень находится на высоте не менее $13{,}2$ метров.

\item 
В физическом эксперименте скорость частицы вычисляется по формуле $v^2 = 10 \cdot 0{,}00025$. Найдите скорость $v$, используя степени десятки.

\item 
Температура нагревателя $T(t) = 800 + 180t - 10t^2$. Прибор портится при температуре $1520$ К. Найдите время (в секундах), через которое прибор необходимо отключить.

\item 
Снаряд вылетает из пусковой установки. Его высота $h(t) = 1{,}5 + 19t - 5t^2$. Найдите общее время (в секундах), в течение которого снаряд находится на высоте не менее $13{,}5$ метров.
\end{enumerate}

\newpage

\SectionTitle{Ответы для самопроверки}

\textcolor{muted}{Ключ-пароль для открытия блока: \(\text{KSENIA-APP-EGE}\).}

\vspace{0.6em}

\renewcommand{\arraystretch}{1.45}
\begin{table}[h]
\centering
\caption{Ответы к тренировочному блоку}
\begin{tabularx}{\linewidth}{@{}cX@{}}
\toprule
\textbf{№} & \textbf{Ответ и краткое пояснение} \\
\midrule
1 & $12$. Уравнение $5P^2 - 90P + 360 = 0 \implies P^2 - 18P + 72 = 0$. Корни $6$ и $12$. Максимальная цена — $12$. \\
2 & $4$. Уравнение $5t^2 - 100t + 320 = 0 \implies t^2 - 20t + 64 = 0$. Корни $4$ и $16$. Отключаем при $t=4$. \\
3 & $1$. Уравнение $5t^2 - 15t + 10 = 0 \implies t^2 - 3t + 2 = 0$. Корни $1$ и $2$. $\Delta t = 2 - 1 = 1$. \\
4 & $0{,}006$. $v = \sqrt{36 \cdot 10^{-6}} = 6 \cdot 10^{-3} = 0{,}006$. \\
5 & $16$. Уравнение $5P^2 - 110P + 480 = 0 \implies P^2 - 22P + 96 = 0$. Корни $6$ и $16$. Максимальная — $16$. \\
6 & $4$. Уравнение $5t^2 - 140t + 480 = 0 \implies t^2 - 28t + 96 = 0$. Корни $4$ и $24$. Отключаем при $t=4$. \\
7 & $1{,}4$. Уравнение $5t^2 - 17t + 12 = 0$. Корни $1$ и $2{,}4$. $\Delta t = 2{,}4 - 1 = 1{,}4$. \\
8 & $0{,}05$. $v^2 = 25 \cdot 10^{-4} \implies v = 5 \cdot 10^{-2} = 0{,}05$. \\
9 & $6$. Уравнение $10t^2 - 180t + 720 = 0 \implies t^2 - 18t + 72 = 0$. Корни $6$ и $12$. Отключаем при $t=6$. \\
10 & $2{,}2$. Уравнение $5t^2 - 19t + 12 = 0$. Корни $0{,}8$ и $3$. $\Delta t = 3 - 0{,}8 = 2{,}2$. \\
\bottomrule
\end{tabularx}
\end{table}

\end{document}